
\medskip
\noindent\SGARefTarget{sga3:iv:lemma:1:7}\textbf{Lemma 1.7.}
\textit{With the notation of \SGARefLink{sga3:iv:lemma:1:4}{1.4}, if \(u\) and \(v\) are effective
epimorphisms and \(v\) is a universal epimorphism, then \(uv\) is an
effective epimorphism.}

To see this, consider the diagram
\[
\begin{array}{ccccc}
  S & \xleftarrow{u} & T & \leftleftarrows & T\times_S T \\
    &                 & \scriptstyle v\,\Big\uparrow
                      &                  & \scriptstyle v\times_Sv\,\Big\uparrow \\[-0.25em]
    &                 & U & \leftleftarrows & U\times_S U \\
    &                 & \substack{\uparrow\\[-0.6ex]\uparrow}
                      & \nearrow & \\
    &                 & U\times_T U & &
\end{array}
\tag{1.7.1}\label{sga3:iv:lemma:1:7:eq:1:7:1:01}
\]
By hypothesis, the first row and the first column are exact; moreover,
by \SGARefLink{sga3:iv:corollary:1:5}{1.5} and \SGARefLink{sga3:iv:lemma:1:6}{1.6}, \(v\times_Sv\) is an epimorphism, since \(v\) is a
universal epimorphism. The conclusion follows by an evident diagram chase.
Indeed, if an element of \(X(U)\) has the same two images in
\(X(U\times_SU)\), then it has, \textit{a fortiori}, the same two images in
\(X(U\times_TU)\), and hence comes from an element of \(X(T)\), since the
first column is exact.

Since the first row is exact, it remains only to verify that the element in
question has the same two images in \(X(T\times_ST)\). Because
\(v\times_Sv\) is an epimorphism, it is enough to verify that their images in
\(X(U\times_SU)\) are the same, which is indeed the case.

\medskip
\noindent\SGARefTarget{sga3:iv:proposition:1:8}\textbf{Proposition 1.8.}
\textit{Consider morphisms \(U\xrightarrow{v}T\xrightarrow{u}S\). If \(u\)
and \(v\) are universal effective epimorphisms, then \(uv\) is a universal
effective epimorphism. If \(uv\) is a universal effective
epimorphism and \(u\) is squarable, then \(u\) is a universal effective
epimorphism.}

The first implication follows at once from \SGARefLink{sga3:iv:lemma:1:7}{1.7}. For the second, consider \SGARefLink{sga3:iv:proposition:1:8:eq:1:8:1:01}{the
following diagram} (of ``bisimplicial'' type):
\[
\begin{array}{ccccc}
  S & \xleftarrow{u} & T & \leftleftarrows & T\times_ST \\
  \scriptstyle uv\,\Big\uparrow
    & & \Big\uparrow & & \Big\uparrow \\[-0.25em]
  U & \longleftarrow & U\times_ST
    & \leftleftarrows & U\times_ST\times_ST \\
  \substack{\uparrow\\[-0.6ex]\uparrow}
    & & \substack{\uparrow\\[-0.6ex]\uparrow}
    & & \substack{\uparrow\\[-0.6ex]\uparrow} \\[-0.25em]
  U\times_SU & \longleftarrow & U\times_SU\times_ST
    & \leftleftarrows & U\times_SU\times_ST\times_ST .
\end{array}
\tag{1.8.1}\label{sga3:iv:proposition:1:8:eq:1:8:1:01}
\]
The first, second, and third columns are exact by the hypothesis that \(uv\)
is a universal effective epimorphism. The second row is exact because
\(U\times_ST\to U\) is an effective epimorphism (it has a section over
\(U\)), and the same is true of the third row, for the same reason. An
evident diagram chase then shows that the first row is exact, that is, that
\(u\) is an effective epimorphism. Since the hypotheses are invariant under
an arbitrary base change \(S'\to S\), it follows that \(u\) is in fact a
universal effective epimorphism.

\medskip
\noindent\SGARefTarget{sga3:iv:corollary:1:9}\textbf{Corollary 1.9.}
\textit{Let \(u:X\to Y\) and \(u':X'\to Y'\) be universal effective
epimorphisms, and suppose that \(Y\times Y'\) exists. Then \(X\times X'\)
exists and}
\[
  u\times u':X\times X'\longrightarrow Y\times Y'
\]
\textit{is a universal effective epimorphism.}

The proof is the same as for \SGARefLink{sga3:iv:corollary:1:5}{1.5}, using the diagram
\[
\begin{array}{ccccc}
   & & Y' & \xleftarrow{u'} & X' \\
   & & \Big\uparrow & & \Big\uparrow \\[-0.25em]
  X & \longleftarrow & X\times Y' & \longleftarrow & X\times X' \\
  \scriptstyle u\,\Big\downarrow
    & & \Big\downarrow & \swarrow_{\scriptstyle u\times u'} & \\[-0.25em]
  Y & \longleftarrow & Y\times Y' & & .
\end{array}
\tag{1.9.1}\label{sga3:iv:corollary:1:9:eq:1:9:1:01}
\]

\medskip
\noindent\SGARefTarget{sga3:iv:corollary:1:10}\textbf{Corollary 1.10.}
\textit{Let \(u:T\to S\) be a squarable morphism, and let \(S'\to S\) be a
base-change morphism that is a universal effective epimorphism. Then \(u\)
is a universal effective epimorphism if and only if}
\[
  u':T'=T\times_SS'\longrightarrow S'
\]
\textit{is one:}
\[
\begin{array}{ccc}
  T & \xleftarrow{v'} & T' \\
  \scriptstyle u\,\Big\downarrow & & \Big\downarrow\,\scriptstyle u' \\[-0.25em]
  S & \xleftarrow{v} & S'.
\end{array}
\tag{1.10.1}\label{sga3:iv:corollary:1:10:eq:1:10:1:01}
\]
Only sufficiency requires proof. If \(u'\) is a universal effective
epimorphism, then so is \(vu'\), by \SGARefLink{sga3:iv:proposition:1:8}{1.8}. Since \(vu'=uv'\), another
application of \SGARefLink{sga3:iv:proposition:1:8}{1.8} shows that \(u\) is a universal effective epimorphism.

\medskip
\noindent\SGARefTarget{sga3:iv:remark:1:11}\textbf{Remark 1.11.}
\textit{The same reasoning shows that in \SGARefLink{sga3:iv:corollary:1:10}{1.10} one may replace ``universal
effective epimorphism'' by ``universal epimorphism,'' by ``universal and
effective epimorphism,'' or simply by ``epimorphism'' (and in the last case
the hypothesis that \(u\) be squarable is plainly unnecessary).}

In the proof of \SGARefLink{sga3:iv:proposition:1:8}{1.8} we used \SGARefLink{sga3:iv:proposition:1:12}{the following result}, which deserves to be stated
explicitly.

\medskip
\noindent\SGARefTarget{sga3:iv:proposition:1:12}\textbf{Proposition 1.12.}
\textit{Let \(u:T\to S\) be a morphism admitting a section. Then \(u\) is an
epimorphism; if \(T\times_ST\) exists, it is an effective epimorphism; and if,
in addition, \(u\) is squarable, it is a universal effective epimorphism.}

The first assertion is contained in \SGARefLink{sga3:iv:lemma:1:4:item:a:01}{1.4(a)}, and the third follows at once
from the second, so it is enough to establish the second. In fact, there is a
stronger conclusion: for every functor
\(F:\mathscr C^{\circ}\to\mathbf{Set}\), not necessarily representable, the
diagram of sets
\[
  F(S)\longrightarrow F(T)\rightrightarrows F(T\times_ST)
\]
is exact. This may be regarded as a special case of the formalism of
\v{C}ech cohomology (in dimension \(0\)!), which we merely recall here.
Assume only that \(T\times_ST\) exists, and set
\[
  \check H^{0}(T/S,F)
  =\operatorname{Ker}\bigl(F(T)\rightrightarrows F(T\times_ST)\bigr).
\]

One may plainly regard \(\check H^{0}(T/S,F)\) as a contravariant functor of
the variable \(T\) in \(\mathscr C_{/S}\): every \(S\)-morphism \(T'\to T\)
defines a map
\[
\tag{+}\label{sga3:iv:proposition:1:12:eq:unnamed:01}
  \check H^{0}(T/S,F)\longrightarrow\check H^{0}(T'/S,F).
\]
Fix \(T,T'\) in \(\mathscr C_{/S}\). A familiar calculation shows that, if
there exists an \(S\)-morphism \(T'\to T\), the corresponding map \SGARefLink{sga3:iv:proposition:1:12:eq:unnamed:01}{(+)} is
independent of the choice of that morphism.\footnote{Editor's note: this is
the following argument, communicated by M.~Demazure. Let
\(f,g:T'\to T\), and let \(\phi:T'\to T\times_ST\) be the morphism with
components \(f\) and \(g\), so that \(p_1\circ\phi=f\) and
\(p_2\circ\phi=g\). Then
\(F(\phi):F(T\times_ST)\to F(T')\) satisfies
\(F(\phi)\circ F(p_1)=F(f)\) and
\(F(\phi)\circ F(p_2)=F(g)\). For every
\(x\in\check H^0(T/S,F)\), however, one has
\(F(p_1)(x)=F(p_2)(x)\). Applying \(F(\phi)\) to both sides gives
\(F(f)(x)=F(g)(x)\), which shows that \(f\) and \(g\) induce the same map.}

Thus \(\check H^{0}(T/S,F)\) may be regarded as a functor on the category
associated with the set \(\operatorname{Ob}(\mathscr C_{/S})\), preordered
by the relation of domination: \(T'\) dominates \(T\) if there exists an
\(S\)-morphism \(T'\to T\). In particular, if \(T\) and \(T'\) are
isomorphic in this latter category, that is, if each dominates the other,
then \SGARefLink{sga3:iv:proposition:1:12:eq:unnamed:01}{(+)} is an isomorphism of sets. This applies in particular when \(T'\)
is the final object of \(\mathscr C_{/S}\), essentially \(S\) itself. The
object \(T\) always dominates \(T'=S\), while the converse holds precisely
when \(T/S\) has a section. This proves \SGARefLink{sga3:iv:proposition:1:12}{1.12} in the stronger form announced.

\medskip
\noindent\SGARefTarget{sga3:iv:remark:1:13}\textbf{Remark 1.13.}
\textit{For various applications, the notions introduced and the results
stated in this expos\'e must be developed more generally for a family of
morphisms \(u_i:T_i\to S\) with the same target, rather than for a single
morphism \(u:T\to S\). Such a family is called \emph{epimorphic} if, for
every object \(X\) of \(\mathscr C\), the corresponding map}
\[
  X(S)\longrightarrow\prod_i X(T_i)
\]
\textit{is injective. One similarly introduces the notion of an effective
epimorphic family and the universal variants of these notions. Whenever
needed below, we shall take for granted that the results of this expos\'e
extend to this more general situation.}

\medskip
\noindent\SGARefTarget{sga3:iv:remark:1:14}\textbf{Remark 1.14.}
\textit{Every morphism that is both a monomorphism and an effective
epimorphism is an isomorphism.}

Indeed, in the notation of \SGARefLink{sga3:iv:definition:1:3}{1.3}, the fact that \(T\to S\) is a monomorphism
implies that the two morphisms
\[
  \operatorname{pr}_1,\operatorname{pr}_2:T\times_ST
  \rightrightarrows T
\]
are equal (and are isomorphisms). But a diagram
\[
  C\overset{v}{\underset{v}{\rightrightarrows}}B\xrightarrow{u}A
\]
is exact only if \(u\) is an isomorphism, as follows immediately from \SGARefLink{sga3:iv:definition:1:3}{the
definition}.\footnote{Editor's note: a monomorphism that is an epimorphism
need not be an isomorphism. For example, in
\(\mathscr C=\mathbf{Sch}\), the morphism
\[
  \operatorname{Spec}(\mathbf F_p)\amalg
  \operatorname{Spec}(\mathbf Z[1/p])
  \longrightarrow\operatorname{Spec}(\mathbf Z)
\]
is a monomorphism and a surjective epimorphism, but is not an isomorphism.}

\SGARefTarget{sga3:iv:section:2}\section*{2. Descent Morphisms}
\addcontentsline{toc}{section}{2. Descent Morphisms}

Recall the following definitions.

\medskip
\noindent\SGARefTarget{sga3:iv:definition:2:1}\textbf{Definition 2.1.}
\textit{Let \(f:S'\to S\) be a morphism such that
\(S''=S'\times_SS'\) exists, and let \(u':X'\to S'\) be an object over
\(S'\). A \emph{gluing datum} on \(X'/S'\), relative to \(f\), is an
\(S''\)-isomorphism}
\[
  c:X''_1\xrightarrow{\sim}X''_2,
\]
\textit{where \(X''_i\), for \(i=1,2\), denotes the inverse image, assumed
to exist, of \(X'/S'\) under the projection
\(\operatorname{pr}_i:S''\to S'\). The gluing datum \(c\) is called a
\emph{descent datum} if it satisfies the ``cocycle condition''}
\[
  \operatorname{pr}_{3,1}^{*}(c)
  =\operatorname{pr}_{3,2}^{*}(c)\,
   \operatorname{pr}_{2,1}^{*}(c),
\]
\textit{where \(\operatorname{pr}_{i,j}\), for
\(1\leq j<i\leq3\), are the canonical projections from}
\[
  S'''=S'\times_SS'\times_SS'
\]
\textit{to \(S''\) (and we now assume that \(S'''\) also exists), where
\(\operatorname{pr}_{i,j}^{*}(c)\) is the inverse image of \(c\), viewed
as an \(S'''\)-morphism from \(X'''_j\) to \(X'''_i\), and where, for each
integer \(k\) from \(1\) to \(3\), \(X'''_k\) denotes the inverse image,
assumed to exist, of \(X'/S'\) under the \(k\)-th projection
\(q_k:S'''\to S'\).}

In the second part of \SGARefLink{sga3:iv:definition:2:1}{the definition} we have therefore used customary
identifications and abuses of notation,\footnote{Editor's note: for example,
we have identified, on the one hand,
\(\operatorname{pr}_{2,1}^{*}(X''_1)
 =\operatorname{pr}_{3,1}^{*}(X''_1)\), and, on the other hand,
\(\operatorname{pr}_{2,1}^{*}(X''_2)
 =\operatorname{pr}_{3,2}^{*}(X''_2)\).} which experience shows to be
harmless, but which must plainly be avoided in a rigorous exposition of
descent theory---a theory that must itself justify, to some extent, these
customary abuses of language. Such a formal exposition (\SGARefLink{sga3:iv:bib:d}{[D]}) was written by
Giraud, with a view to justifying and making precise \SGARefLink{sga3:iv:bib:sga:1}{SGA~1}, Expos\'e~VII,
which was never written. For a detailed account of the faithfully flat
descent results used constantly in this Seminar, one may consult \SGARefLink{sga3:iv:bib:sga:1}{SGA~1},
Expos\'e~VIII.

Continue to let \(f:S'\to S\) be a morphism for which
\(S''=S'\times_SS'\) exists, and let \(X\) be an object over \(S\) such
that \(X'=X\times_SS'\) and \(X''=X\times_SS''\) exist. The inverse images
of \(X'\) under \(\operatorname{pr}_i\), for \(i=1,2\), then exist and are
canonically isomorphic; consequently, \(X'/S'\) carries a canonical gluing
datum relative to \(f\). When \(S'''\) and \(X'''=X\times_SS'''\) exist,
this is in fact a descent datum. If \(Y\) is another object over \(S\)
satisfying the same conditions as \(X/S\), then for every \(S\)-morphism
\(X\to Y\), the corresponding \(S'\)-morphism \(X'\to Y'\) is compatible
with the canonical gluing data on \(X'\) and \(Y'\). In particular, if
\(S'\to S\) is squarable, then
\[
  X\longmapsto X'=X\times_SS'
\]
is a functor from \(\mathscr C_{/S}\) to the ``category of objects over
\(S'\) equipped with a descent datum relative to \(f\)''---a category whose
definition is left to the reader and which is a full subcategory of the
``category of objects over \(S'\) equipped with a gluing datum relative to
\(f\).''

With this understood, we make \SGARefLink{sga3:iv:definition:2:2}{the following definition}.

\medskip
\noindent\SGARefTarget{sga3:iv:definition:2:2}\textbf{Definition 2.2.}
\textit{A morphism \(f:S'\to S\) is called a \emph{descent morphism}
(respectively, an \emph{effective descent morphism}) if \(f\) is squarable
(that is, \(X\times_SS'\) exists for every \(X\to S\)) and the preceding
functor}
\[
  X\longmapsto X'=X\times_SS'
\]
\textit{from the category \(\mathscr C_{/S}\) of objects over \(S\) to the
category of objects over \(S'\) equipped with a descent datum relative to
\(f\) is fully faithful (respectively, an equivalence of categories).}

The first of these notions can be expressed using gluing data alone, and
therefore without invoking the triple fiber product \(S'''\): \(f\) is a
descent morphism precisely when it is squarable and \(X\mapsto X'\) is a
fully faithful functor from \(\mathscr C_{/S}\) to the category of objects
over \(S'\) equipped with a gluing datum relative to \(f\). Written out, this
means that for any two objects \(X,Y\) over \(S\), \SGARefLink{sga3:iv:definition:2:2:eq:x:01}{the following diagram} of
sets is exact:
\[
\tag{x}\label{sga3:iv:definition:2:2:eq:x:01}
  \operatorname{Hom}_S(X,Y)
  \longrightarrow \operatorname{Hom}_{S'}(X',Y')
  \overset{p_1}{\underset{p_2}{\rightrightarrows}}
  \operatorname{Hom}_{S''}(X'',Y'').
\]

Here \(p_i(u')\) denotes the inverse image of
\(u'\in\operatorname{Hom}_{S'}(X',Y')\) under the projection
\(\operatorname{pr}_i:S''\to S'\), for \(i=1,2\). Indeed, by definition,
the kernel of the pair \((p_1,p_2)\) is precisely the set of
\(S'\)-morphisms \(X'\to Y'\) compatible with the canonical gluing data.

By the definition of the inverse images \(Y'\) and \(Y''\), there are
canonical bijections
\[
  \operatorname{Hom}_{S'}(X',Y')\simeq\operatorname{Hom}_S(X',Y),
  \qquad
  \operatorname{Hom}_{S''}(X'',Y'')\simeq\operatorname{Hom}_S(X'',Y).
\]
Thus exactness of \SGARefLink{sga3:iv:definition:2:2:eq:x:01}{(x)} is equivalent to exactness of
\[
\tag{xx}\label{sga3:iv:definition:2:2:eq:xx:01}
  \operatorname{Hom}_S(X,Y)
  \longrightarrow\operatorname{Hom}_S(X',Y)
  \rightrightarrows\operatorname{Hom}_S(X'',Y),
\]
which is obtained by applying \(\operatorname{Hom}_S(-,Y)\) to the diagram
in \(\mathscr C_{/S}\)
\[
\tag{xxx}\label{sga3:iv:definition:2:2:eq:xxx:01}
  X''\rightrightarrows X'\longrightarrow X,
\]
itself obtained from
\[
  S''\rightrightarrows S'\longrightarrow S
\]
by the base change \(X\to S\). In view of \SGARefLink{sga3:iv:lemma:1:6}{1.6}, this proves the first part of
\SGARefLink{sga3:iv:proposition:2:3}{the following proposition}.

\medskip
\noindent\SGARefTarget{sga3:iv:proposition:2:3}\textbf{Proposition 2.3.}
\textit{Let \(f:S'\to S\) be a morphism. If \(f\) is a universal effective
epimorphism, then it is a descent morphism. Conversely, every descent
morphism is a universal effective epimorphism if \(S\) is a squarable object
of \(\mathscr C\), that is, if its product with every object of
\(\mathscr C\) exists.}

It remains to prove the converse. If \(f\) is a descent morphism, we must
show that it is a universal effective epimorphism: for every morphism
\(X\to S\), diagram \SGARefLink{sga3:iv:definition:2:2:eq:xxx:01}{(xxx)} must be exact, or equivalently, its image under
\(\operatorname{Hom}(-,Z)\) must be an exact diagram of sets for every
object \(Z\) of \(\mathscr C\). By hypothesis, \(Z\times S\) exists. Let
\(Y\) be the object of \(\mathscr C_{/S}\) that it defines. The image of
\SGARefLink{sga3:iv:definition:2:2:eq:xxx:01}{(xxx)} under \(\operatorname{Hom}(-,Z)\) is then isomorphic to its image
under \(\operatorname{Hom}_S(-,Y)\), and the latter is exact by the
hypothesis on \(f\).

The results concerning descent morphisms may therefore be applied to
universal effective epimorphisms. For example:

\medskip
\noindent\SGARefTarget{sga3:iv:proposition:2:4}\textbf{Proposition 2.4.}
\textit{Let \(f:S'\to S\) be a descent morphism (for example, a universal
effective epimorphism). Then:}
\begin{enumerate}
  \item[(a)]\SGARefTarget{sga3:iv:proposition:2:4:item:a:01} \textit{For every \(S\)-morphism \(u:X\to Y\), the morphism
    \(u\) is an isomorphism (respectively, a monomorphism) if and only if
    \(u':X'\to Y'\) is one.}
  \item[(b)]\SGARefTarget{sga3:iv:proposition:2:4:item:b:01} \textit{Let \(X,Y\) be two subobjects of \(S\), and let
    \(X',Y'\) be their inverse images in \(S'\). Then \(X\) is contained in
    \(Y\) (respectively, is equal to \(Y\)) if and only if the same is true
    of \(X'\) and \(Y'\).}
\end{enumerate}

For \SGARefLink{sga3:iv:proposition:2:4:item:a:01}{(a)}, \SGARefLink{sga3:iv:definition:2:2}{the definition} shows that if \(u'\) is an isomorphism in the
category of objects with gluing datum, then \(u\) is an isomorphism. Every
isomorphism of objects over \(S'\) that is compatible with gluing data is,
however, an isomorphism of objects with gluing datum: its inverse is also
compatible with the gluing data. For \SGARefLink{sga3:iv:proposition:2:4:item:b:01}{(b)}, it remains to show that if \(X'\)
is contained in \(Y'\), that is, if there exists an \(S'\)-morphism
\(X'\to Y'\), then the same is true of \(X,Y\) over \(S\). Since
\(Y'\to S'\), and hence \(Y''\to S''\), is a monomorphism, the morphism
\(X'\to Y'\) is automatically compatible with the gluing data and therefore
comes from an \(S\)-morphism \(X\to Y\).

The proof applies more generally to two objects \(X,Y\) over \(S\), with
\(Y\to S\) a monomorphism, when one asks whether \(X\to S\) factors through
\(Y\): it is enough that \(X'\to S'\) factor through \(Y'\).

\medskip
\noindent\SGARefTarget{sga3:iv:corollary:2:5}\textbf{Corollary 2.5.}
\textit{Let \(f:S'\to S\) be a universal effective epimorphism and
\(g:S\to T\) a morphism such that \(S\times_TS\) exists. Suppose that
\(S''=S'\times_SS'\) is also a fiber product of \(S'\) with itself over
\(T\), that is, that the canonical morphism}
\[
  S'\times_SS'\xrightarrow{\sim}S'\times_TS'
\]
\textit{is an isomorphism. Then \(g:S\to T\) is a monomorphism (and the
converse is of course true).}

Indeed, consider the cartesian diagram
\[
\begin{array}{ccc}
  S'\times_SS' & \xrightarrow{\sim} & S'\times_TS' \\
  \Big\downarrow & & \Big\downarrow\,\scriptstyle f \\[-0.25em]
  S & \longrightarrow & S\times_TS ,
\end{array}
\tag{2.5.1}\label{sga3:iv:corollary:2:5:eq:2:5:1:01}
\]
where the second horizontal arrow is the diagonal morphism. By \SGARefLink{sga3:iv:corollary:1:9}{1.9}, the
second vertical arrow \(f\) is a universal effective epimorphism; by
hypothesis, the first horizontal arrow is an isomorphism. It follows from
\SGARefLink{sga3:iv:proposition:2:4:item:a:01}{2.4(a)}, or equally from \SGARefLink{sga3:iv:proposition:2:4:item:b:01}{2.4(b)},\footnote{Editor's note: applied to
\(f:S'\times_TS'\to S\times_TS=Y\).} that
\(S\to S\times_TS\) is also an isomorphism, which says precisely that
\(g:S\to T\) is a monomorphism.

\medskip
\noindent\SGARefTarget{sga3:iv:remark:2:6}\textbf{Remark 2.6.}
\textit{The notions introduced in \SGARefLink{sga3:iv:definition:2:1}{2.1} in terms of morphisms between certain
inverse limits have an evident expression in terms of the contravariant
functors defined by the objects \(S,S',X'\) under consideration. Subject to
the existence of the fiber products considered in \SGARefLink{sga3:iv:definition:2:1}{2.1}, there is a one-to-one
correspondence between gluing data (respectively, descent data) on \(X'/S'\)
relative to \(S'\to S\), and gluing data (respectively, descent data) for the
corresponding objects of}
\[
  \widehat{\mathscr C}
  =\operatorname{Hom}(\mathscr C^{\circ},\mathbf{Set}).
\]
\textit{This makes it possible, if desired, to extend these notions to the
case in which no existence hypothesis is imposed on fiber products in
\(\mathscr C\).}

\medskip
\noindent\SGARefTarget{sga3:iv:remark:2:7}\textbf{Remark 2.7.}
\textit{The notions introduced in \SGARefLink{sga3:iv:section:2}{this section} generalize to families of
morphisms. They can also be presented more intrinsically by means of the
notion of a sieve (\SGARefLink{sga3:iv:subsection:4:1}{4.1}). For these matters the reader is referred to \SGARefLink{sga3:iv:bib:d}{[D]}.}
